\chapter*{Заключение} \label{ch-conclusion}
\addcontentsline{toc}{chapter}{Заключение}

В ходе практической работы выполнена разработка BDD-тестирования для консольного приложения на Node.js, решающего задачу поиска палиндромных подстрок алгоритмом Манакера. Для приложения подготовлены три Gherkin-истории: анализ палиндромов, JSON-ввод-вывод и нагрузочное сравнение двух алгоритмов. Сценарии содержат пользовательские истории, \texttt{Rule}, \texttt{Background}, \texttt{Scenario Outline}, таблицы входных данных, а также шаги \texttt{Given}, \texttt{When}, \texttt{Then}, \texttt{And} и \texttt{But}.

В качестве BDD-фреймворка выбран Cucumber.js, так как он напрямую поддерживает Gherkin и хорошо подходит для проекта на JavaScript. Для сценариев реализованы Step Definitions, которые вызывают реальные модули приложения: сервис анализа палиндромов, JSON-чтение, запись результата и скрипты нагрузочного эксперимента.

Экспериментальная часть подтвердила теоретические оценки алгоритмов. Алгоритм Манакера показал линейное поведение и на всех наборах выполнялся за время порядка единиц миллисекунд. Эталонное расширение от центра оказалось чувствительным к структуре строки: на наборе \texttt{UNI\_16000\_A} оно заняло 270.883 мс против 1.419 мс у Манакера. Наблюдаемое ускорение составило примерно 190.83 раза, что соответствует различию между $O(n)$ и $O(n^2)$ в худшем случае.

BDD и модульные тесты решают разные задачи. Unit-тесты быстрее и точнее проверяют внутренние функции, массивы радиусов и граничные условия. BDD-сценарии медленнее, но лучше выражают пользовательские требования и служат исполняемой документацией. Поэтому наиболее практичным оказался комбинированный подход: unit-тесты сохраняются как техническая проверка, а BDD-сценарии фиксируют ожидаемое поведение приложения.

Таким образом, цель работы достигнута: приложение доработано под BDD-подход, подготовлены автотесты и датасеты, проведено сравнение алгоритмов, а отчёт оформлен по структуре задания.
